%======================================================================
% STAR Test Report
% Capstone Deliverable -- Spring 2026
% Self-contained single-file LaTeX source, 7-bit ASCII only.
% Compile: pdflatex STAR_Test_Report.tex (run twice)
%======================================================================
\documentclass[11pt,letterpaper]{article}

%----------------------------------------------------------------------
% Packages (paired preamble with the STAR SDD and SSUM)
%----------------------------------------------------------------------
\usepackage[margin=1in]{geometry}
\usepackage[T1]{fontenc}
\usepackage{lmodern}
\usepackage{xcolor}
\usepackage{graphicx}
\usepackage{booktabs}
\usepackage{array}
\usepackage{tabularx}
\usepackage{longtable}
\usepackage{enumitem}
\usepackage{parskip}
\usepackage{float}
\usepackage{caption}
\usepackage{fancyhdr}
\usepackage{titlesec}
\usepackage{amsmath}
\usepackage{amssymb}
\usepackage{siunitx}
\usepackage{listings}
\usepackage{tikz}
\usetikzlibrary{positioning, arrows.meta, shapes.geometric}
\usepackage[colorlinks=true,
            linkcolor=blue!55!black,
            citecolor=blue!55!black,
            urlcolor=blue!55!black,
            pdftitle={STAR Test Report},
            pdfauthor={Cesar Magana (Team Locked Inc)},
            pdfsubject={STAR Capstone Test Report},
            pdfkeywords={STAR, testing, Unity, Vitest, pytest, CI, coverage,
                         IEEE 29119}]{hyperref}

%----------------------------------------------------------------------
% Status-label hyperlink macro (paired with SDD and SSUM)
% Use as: \statuslabel{STRETCH} etc. All occurrences hyperlink back
% to the legend at sec:status-legend in the Introduction.
%----------------------------------------------------------------------
\newcommand{\statuslabel}[1]{\hyperref[sec:status-legend]{[#1]}}

\pagestyle{fancy}
\fancyhf{}
\fancyhead[L]{\small STAR Test Report}
\fancyhead[R]{\small Spring 2026}
\fancyfoot[C]{\thepage}
\renewcommand{\headrulewidth}{0.4pt}
\renewcommand{\footrulewidth}{0.0pt}

\titleformat{\section}{\Large\bfseries}{\thesection}{0.75em}{}
\titleformat{\subsection}{\large\bfseries}{\thesubsection}{0.6em}{}
\titleformat{\subsubsection}{\normalsize\bfseries}{\thesubsubsection}{0.5em}{}
\titlespacing*{\section}{0pt}{14pt}{6pt}
\titlespacing*{\subsection}{0pt}{10pt}{4pt}

\definecolor{codegray}{gray}{0.45}
\definecolor{boxblue}{RGB}{230,240,255}
\definecolor{boxborder}{RGB}{70,110,170}
\definecolor{passgreen}{RGB}{40,140,60}
\definecolor{failred}{RGB}{175,40,40}

\lstdefinestyle{shellcmd}{
  frame=single,
  basicstyle=\small\ttfamily,
  breaklines=true,
  numbers=none,
  captionpos=b,
  keepspaces=true,
  showstringspaces=false,
  xleftmargin=0.8em,
  framexleftmargin=0.4em
}

\sisetup{per-mode=symbol, detect-weight=true, detect-family=true}

\tikzset{
  sddbox/.style = {rectangle, draw=boxborder, fill=boxblue, rounded corners,
                   minimum width=2.4cm, minimum height=0.9cm, align=center,
                   font=\small},
  sddlabel/.style = {font=\footnotesize\itshape},
  sddarrow/.style = {-Latex, thick, draw=boxborder}
}

\newcommand{\pass}{{\color{passgreen}\textbf{PASS}}}
\newcommand{\fail}{{\color{failred}\textbf{FAIL}}}

%======================================================================
% Title Page
%======================================================================
\begin{document}
\thispagestyle{empty}

\begin{center}
\vspace*{1.0in}
{\Huge \bfseries Test Report}\\[0.35in]
{\LARGE STAR: Spatial Topography Accessibility Robot}\\[0.20in]
{\large Unit, Integration, and Hardware-in-the-Loop Test Results}\\[0.8in]

\rule{4.5in}{0.8pt}\\[0.3in]

{\large \textbf{Team Locked Inc}}\\[0.15in]
\begin{tabular}{r l}
Cesar Magana        & Software, Project Manager \\
Brighton Sikarskie  & Software, Embedded Development, Schematic Design \\
Jeremie Hockey      & PCB Design and Layout \\
Michael Norton      & Mechanical Design \\
Shawn Ciaciura      & Schematic Review \\
Jared Bartz         & Schematic Review \\
\end{tabular}\\[0.4in]

{\large \textbf{Course}}\\[0.1in]
{\normalsize ESET 420 -- Senior Design II -- Dr. Lee Hudson}\\[0.4in]

{\large \textbf{Institution}}\\[0.1in]
{\normalsize Texas A\&M University}\\
{\normalsize Department of Electronic Systems Engineering Technology (ESET)}\\
{\normalsize ESET Senior Capstone, Spring 2026}\\[0.4in]

{\large \textbf{Advisors and Sponsors}}\\[0.1in]
{\normalsize Dr. Logan Porter -- Sponsor and Technical Advisor}\\[0.6in]

\rule{4.5in}{0.8pt}\\[0.2in]
{\large April 28, 2026}\\[0.15in]
{\small Version 1.0 -- Capstone Submission}

\end{center}
\clearpage

%======================================================================
% Front Matter
%======================================================================
\pagenumbering{roman}
\setcounter{page}{1}
\tableofcontents
\clearpage
\listoffigures
\listoftables
\clearpage

\pagenumbering{arabic}
\setcounter{page}{1}

%======================================================================
% 1. Introduction
%======================================================================
\section{Introduction}
\label{sec:intro}

\subsection{Purpose}
This report documents the verification and validation effort for the
STAR software. It enumerates every automated test suite, identifies
the subsystem each suite exercises, reports pass/fail status and
measured coverage where available, and records the hardware
bring-up tests performed against real silicon. Its structure is
informed by ISO/IEC/IEEE 29119-3:2021 (Software and systems
engineering -- Software testing -- Part 3: Test documentation) and
IEEE 829-2008 legacy practice (superseded by 29119-3), adapted to
the multi-language, multi-target STAR codebase.

\subsection{Scope}
The report covers eight software subsystems: the RX72N real-time
firmware, the Go gateway, the ROS2 Jazzy autonomy stack, the
TypeScript operator UI, the Python compliance engine, the Protocol
Buffers schema repository, the virtual-RX72N simulator, and the CI/CD
automation that drives the whole pipeline. Physical PCB and
mechanical verification are out of scope for this report.

\subsection{Companion Documents and Source-of-Truth Repository}
The STAR project is open-source. The full repository, including
every CI workflow, test harness, and coverage artifact referenced
in this report, lives at
\url{https://github.com/Locked-Inc/STAR}. The two companion
documents below cover, respectively, the design that this report
verifies and the operational procedures that follow from it:
\begin{itemize}[leftmargin=1.4em,itemsep=2pt]
\item STAR Software Description Document (SDD):
      \href{https://github.com/Locked-Inc/STAR/tree/main/final_docs/software_description_document}{\texttt{final\_docs/software\_description\_document/}}.
      Authoritative for any design decision this report mentions
      without elaboration.
\item STAR System and Software User Manual (SSUM):
      \href{https://github.com/Locked-Inc/STAR/tree/main/final_docs/system_software_user_manual}{\texttt{final\_docs/system\_software\_user\_manual/}}.
      Operator-facing procedures.
\end{itemize}
Concrete artifacts cited later in this report -- coverage outputs,
JUnit XML, CI workflow YAMLs, test source files -- can each be
opened directly on GitHub at the path given.

\subsection{Status Label Legend}
\phantomsection
\label{sec:status-legend}
A handful of test scopes in this report are tagged with
\texttt{[STRETCH]} when they were intentionally de-prioritized
during the semester. The same three-tier taxonomy used in the SDD
and SSUM applies; bracket tags elsewhere in this report hyperlink
back to this legend.
\begin{itemize}[leftmargin=1.4em,itemsep=2pt]
\item \textbf{[IMPLEMENTED]} -- complete and exercised in daily use.
\item \textbf{[STRETCH]} -- partial implementation with scaffolding;
      fit for demonstration but not yet complete.
\item \textbf{[ARCHITECTED]} -- designed and specified but not coded.
\end{itemize}

\subsection{Test Strategy at a Glance}
STAR applies a four-tier strategy:
\begin{enumerate}[leftmargin=1.4em,itemsep=2pt]
\item \textbf{Unit tests} run on host with native compilers and
      mock-injected hardware (Unity for firmware, Go's \texttt{testing}
      for the gateway, \texttt{pytest} for compliance, Vitest for UI
      and Protobuf TypeScript, C Unity for nanopb).
\item \textbf{Integration tests} exercise multi-component flows via
      an in-process simulator (\texttt{virtual\_rx72n}) and ROS2
      launch tests.
\item \textbf{Hardware-in-the-loop (HIL) tests} execute nightly on a
      dedicated fixture and against bench firmware-bring-up dirs.
\item \textbf{CI-enforced quality gates} block merges on coverage
      regression, lint failure, protobuf breaking changes, and 7-bit
      ASCII violations.
\end{enumerate}

%======================================================================
% 2. Test Environment
%======================================================================
\section{Test Environment}
\label{sec:env}

\subsection{Host Toolchain}
Unit-test runs target x86\_64 Linux (Ubuntu 24.04) on CI and on
developer workstations. Toolchains include GNU Make, CMake, the
GNURX cross-compiler for firmware host-test fallback, Go 1.26,
Node 22 with Vitest, Python 3.12, and ROS2 Jazzy.

\subsection{Target Hardware}
HIL and bench tests run against: Renesas RX72N (4 MB Flash, 512 KB
SRAM) on the STAR custom PCB, four DRV8263H H-bridge motor drivers,
quadrature encoders (341 PPR), Bosch BNO055 IMU, BMP280 barometer,
HC-SR04 ultrasonic sensors, RPLiDAR C1 2D LiDAR, and a Raspberry
Pi 5 host.

\subsection{Test Frameworks}
Table~\ref{tab:frameworks} maps each subsystem to its test framework.

\begin{table}[H]
\centering
\caption{Test frameworks in use by subsystem.}
\label{tab:frameworks}
\small
\begin{tabularx}{\linewidth}{l l X}
\toprule
\textbf{Subsystem} & \textbf{Framework} & \textbf{Notes} \\
\midrule
star-rx72n-firmware   & Unity v2.6.1        & CMake FetchContent; 41 mock \texttt{.c} files (49 mock \texttt{.h}) including ThreadX \\
star-gateway          & Go \texttt{testing} & With \texttt{-race} and \texttt{-coverprofile} \\
star-ros2 (C++)       & ament\_cmake\_gtest & colcon test integration \\
star-ros2 (Python)    & ament\_python / pytest & launch\_test for simulation \\
star-ui               & Vitest              & jsdom environment, testing-library \\
compliance-engine     & pytest              & Coverage via \texttt{pytest-cov} \\
star-proto (Go)       & Go \texttt{testing} & Serialization round-trips \\
star-proto (TS)       & Vitest              & Serialization + conformance \\
star-proto (nanopb)   & Unity               & C serialization tests \\
\bottomrule
\end{tabularx}
\end{table}

%======================================================================
% 3. Firmware Unit Tests
%======================================================================
\section{Firmware Unit Tests (RX72N)}
\label{sec:fw}

\subsection{Scope and Inventory}
Unity-based tests live under
\texttt{star-rx72n-firmware/tests/}. Sixty test source files exercise
the firmware libraries in isolation, using function-pointer dependency
injection to replace hardware with mocks. Tests are grouped by
concern: cores (\texttt{rx\_crc8}, \texttt{rx\_crc32},
\texttt{rx\_frame}, \texttt{rx\_fec}, \texttt{rx\_harq},
\texttt{rx\_pid}, \texttt{rx\_error\_handler},
\texttt{rx\_pin\_validator}, \texttt{rx\_register\_guard});
hardware abstraction
(\texttt{gpio\_hal}, \texttt{uart\_hal}, \texttt{rx\_gptw\_staggered},
\texttt{rx\_mpc}, \texttt{rx\_usb}, \texttt{rx\_usb\_multiport},
\texttt{rx\_usb\_comm});
communication (\texttt{rx\_comm\_manager}, \texttt{rx\_spi\_comm},
\texttt{rx\_bus\_manager}, \texttt{rx\_bus\_i2c},
\texttt{rx\_bus\_gpio}, \texttt{rx\_bus\_onewire},
\texttt{rx\_bus\_adc}, \texttt{rx\_nanopb});
sensors (\texttt{rx\_hcsr04}, \texttt{rx\_obstacle\_detect},
\texttt{rx\_ds18b20}, \texttt{rx\_encoder},
\texttt{rx\_encoder\_tpu}, \texttt{rx\_bmp280},
\texttt{rx\_bno055}); motor control (\texttt{rx\_motor}); and the
seven ThreadX task entry points.

\subsection{Results}
Table~\ref{tab:fw-results} summarizes the current firmware test
suite as wired into \texttt{star-rx72n-firmware/tests/CMakeLists.txt}
(60 \texttt{test\_*.c} files, all registered as ctest entries). The
e2-studio JUnit snapshot at
\texttt{e2-studio-star-rx72n-firmware/tests/build/test-results.xml}
captured 46 of these suites on 2026-03-01; the remaining 14 were
added on the parallel CMake host tree later and are exercised by
the \texttt{firmware-unit-tests.yml} CI gate, which reports the
authoritative current count.

\begin{table}[H]
\centering
\caption{Firmware unit test inventory and most-recent CI results.}
\label{tab:fw-results}
\small
\begin{tabular}{l r}
\toprule
\textbf{Metric} & \textbf{Value} \\
\midrule
Test source files (\texttt{tests/test\_*.c})          & 60 \\
ctest entries registered                              & 60 \\
Failures                                              & 0 \\
Errors                                                & 0 \\
Skipped                                               & 0 \\
Overall status                                        & \pass{} \\
Wall-clock duration (per case)                        & \qtyrange{15}{25}{\milli\second} \\
e2-studio JUnit snapshot (2026-03-01) suites covered  & 46 \\
\bottomrule
\end{tabular}
\end{table}

\subsection{Coverage Gate}
The firmware CI workflow \texttt{.github/workflows/firmware-unit-tests.yml}
generates an lcov HTML coverage report and applies a strict gate: \textbf{100
percent line, branch, and function coverage on
\texttt{libs/}} (generated code and ThreadX task bodies are excluded).
A regression below this threshold blocks merge.

\subsection{Notable Test Categories}
\begin{itemize}[leftmargin=1.4em,itemsep=2pt]
\item \textbf{PID correctness} (\texttt{test\_rx\_pid}): checks the
      discrete velocity-form update against a scalar reference
      implementation and verifies anti-windup clamping behavior.
\item \textbf{CRC-32 conformance} (\texttt{test\_rx\_crc32}): agrees
      with the Go gateway's CRC-32 on a shared fixture vector,
      guaranteeing cross-language byte identity.
\item \textbf{Frame go-compat} (\texttt{test\_frame\_go\_compat}):
      encodes a set of payloads in firmware code and checks that
      the Go frame decoder accepts them bit-exactly.
\item \textbf{FEC and HARQ} (\texttt{test\_rx\_fec},
      \texttt{test\_rx\_harq}): validates the rate-1/2 K=7
      convolutional encoder against the Viterbi decoder and
      verifies Chase Combining soft-bit accumulation across
      simulated retransmissions.
\end{itemize}

\subsection{Hardware Bring-Up Tests}
In addition to unit tests, stand-alone firmware trees at the root of
\texttt{star-rx72n-firmware/} provide on-target bring-up programs used
during PCB validation. Table~\ref{tab:bringup} lists each test and its
current status as exercised on the bench.

\begin{table}[H]
\centering
\caption{Hardware bring-up test directories.}
\label{tab:bringup}
\small
\begin{tabularx}{\linewidth}{l X l}
\toprule
\textbf{Directory} & \textbf{What it verifies} & \textbf{Status} \\
\midrule
\texttt{motor\_spin\_test/}   & Open-loop PWM drive on all four motors              & \pass{} \\
\texttt{pwm\_test/}           & \qty{20}{\kilo\hertz} PWM waveform + dead-time on GPTW & \pass{} \\
\texttt{encoder\_test/}       & TPU quadrature counts and direction sign             & \pass{} \\
\texttt{gpio\_test/}          & 100+ GPIO toggles via external AD2 fixture           & \pass{} \\
\texttt{uart\_test/}          & SCI loopback against debug USB-UART bridge          & \pass{} \\
\texttt{imu\_test/}           & BNO055 orientation output over I2C                  & \pass{} \\
\texttt{sonar\_test/}         & HC-SR04 trigger/echo and range computation          & \pass{} \\
\texttt{usb\_test/}           & USB CDC enumeration and byte echo                   & \pass{} \\
\bottomrule
\end{tabularx}
\end{table}

\subsection{Manual PCB Bring-Up and Validation}
The on-target firmware programs in Table~\ref{tab:bringup} run after
the board is electrically validated. Before any firmware was flashed,
the assembled custom motherboard PCB went through a formal bench
validation pass to verify the layout, the power tree, the motor
driver wiring, and the digital fabric. Table~\ref{tab:pcbvalid}
records each manual check, the instrument used, the pass criterion,
and the outcome on the production board. This is the first formal
write-up of these checks; the bench notes were captured in real
time on the engineering laptop and are summarized below as the
authoritative record.

\begin{table}[H]
\centering
\caption{Manual PCB validation tests on the production STAR motherboard.}
\label{tab:pcbvalid}
\small
\setlength{\tabcolsep}{5pt}
\renewcommand{\arraystretch}{1.15}
\begin{tabularx}{\linewidth}{@{}%
  >{\raggedright\arraybackslash}p{3.0cm}%
  >{\raggedright\arraybackslash}p{3.0cm}%
  >{\raggedright\arraybackslash}X%
  >{\centering\arraybackslash}p{1.0cm}@{}}
\toprule
\textbf{Test} & \textbf{Instrument} & \textbf{Method and pass criterion} & \textbf{Status} \\
\midrule
Power-rail continuity
  & DMM (continuity + DCV)
  & \qty{3.3}{\volt}, \qty{5.0}{\volt}, and \qty{6.0}{\volt} buck outputs
    buzz to every downstream load pin; no short to GND; each rail measures
    within $\pm\qty{5}{\percent}$ of nominal at idle.
  & \pass{} \\[2pt]
GND-net continuity
  & DMM (continuity)
  & All ground vias, copper pours, connector shells, and IC GND pins tie
    to a single net per the KiCad schematic.
  & \pass{} \\[2pt]
Motor-driver output continuity
  & DMM (continuity)
  & DRV8263H \texttt{OUT1} and \texttt{OUT2} on each of the four H-bridge
    channels buzz to the corresponding motor-connector pin pair; no
    cross-channel short.
  & \pass{} \\[2pt]
Net-vs-schematic verification
  & DMM (continuity)
  & Each signal net buzzed end-to-end against the KiCad netlist export;
    every observed connection matched, no spurious connections.
  & \pass{} \\[2pt]
Free-spin motor test
  & DMM (DC current) + \texttt{motor\_spin\_test/}
  & Chassis on blocks (wheels free); PWM ramped 0--100\,\% on each of four
    motors; commanded direction, smooth spin-up, and steady-state no-load
    current within driver envelope.
  & \pass{} \\[2pt]
Loaded-chassis motor test
  & DMM (DC current) + \texttt{motor\_spin\_test/}
  & Chassis on the floor under its own weight; straight-line drive plus
    turn-in-place; all four motors deliver torque, loaded current within
    DRV8263H envelope, no thermal shutdown over a multi-minute run.
  & \pass{} \\[2pt]
PWM-waveform probe
  & Oscilloscope
  & DRV8263H \texttt{IN1}/\texttt{IN2} measured at \qty{20}{\kilo\hertz}
    fundamental, commanded duty cycle, with configured dead-time visible
    at edge transitions.
  & \pass{} \\[2pt]
Encoder-quadrature probe
  & AD2 / oscilloscope
  & A/B phase pair captured while a wheel was rotated by hand;
    \qty{90}{\degree} phase relationship and correct increment/decrement
    on direction reversal on all four encoders.
  & \pass{} \\[2pt]
Comm-line probe (I2C, UART, SPI)
  & AD2 / oscilloscope
  & Idle state, clean clocking at expected baud / SCLK rate, no glitches;
    sampled during normal firmware operation on each interface.
  & \pass{} \\
\bottomrule
\end{tabularx}
\end{table}

These checks were performed on the production-rev STAR motherboard
prior to integrating the chassis. Every test passed first-pass; no
PCB rework was required before the firmware bring-up tree in
Table~\ref{tab:bringup} ran on the same board.

%======================================================================
% 4. Gateway Unit and Integration Tests
%======================================================================
\section{Gateway Tests (Go)}
\label{sec:gw}

\subsection{Scope and Inventory}
Forty-six \texttt{*\_test.go} files across fourteen packages exercise
the gateway's protocol stack and service layer. Key package-level
test groups are shown in Table~\ref{tab:gw-packages}.

\begin{table}[H]
\centering
\caption{Gateway test packages.}
\label{tab:gw-packages}
\small
\begin{tabularx}{\linewidth}{l X}
\toprule
\textbf{Package} & \textbf{Focus} \\
\midrule
\texttt{internal/frame}      & Framing encode/decode, stream decoder, fuzz \\
\texttt{internal/fec}        & Rate-1/2 K=7 convolutional encode/Viterbi decode \\
\texttt{internal/harq}       & Chase Combining, mock transport integration \\
\texttt{internal/dispatcher} & Message routing, fuzzing \\
\texttt{internal/manager}    & Session, heartbeat, health, hot-plug detection \\
\texttt{internal/link}       & CDC edge cases, SPI framing \\
\texttt{internal/transport}  & CDC and SPI mocks \\
\texttt{internal/controller} & Drive control, handler lifecycle \\
\texttt{internal/service}    & gRPC services: motor, telemetry, firmware,
                               configuration, gateway control \\
\texttt{internal/server}     & gRPC, HTTP, config wiring \\
\texttt{internal/ws}         & WebSocket hub and handler \\
\texttt{cmd/virtual\_rx72n}  & In-process RX72N simulator \\
\texttt{test/e2e}            & End-to-end HIL test harness \\
\bottomrule
\end{tabularx}
\end{table}

\subsection{Results and Coverage}
Gateway tests run under \texttt{go test -race
-coverprofile=coverage.out -covermode=atomic ./...}. The race
detector is always on. Coverage is enforced at
\texttt{MIN\_COVERAGE=65\%} by
\texttt{.github/workflows/gateway.yml}.

\begin{table}[H]
\centering
\caption{Gateway coverage summary from committed
\texttt{star-gateway/coverage.out}.}
\label{tab:gw-cov}
\small
\begin{tabular}{l r}
\toprule
\textbf{Metric} & \textbf{Value} \\
\midrule
Packages with tests          & 14 \\
Total test files             & 46 \\
Total statement coverage     & 74.0\% \\
CI gate                      & 65.0\% \\
Status vs.\ gate             & \pass{} (+9.0 points above gate) \\
Race detector                & enabled \\
\bottomrule
\end{tabular}
\end{table}

Sub-package coverage spot-checks: \texttt{internal/fec} and
\texttt{internal/frame/decoder.go} measure at 100.0 percent. In
the \texttt{ws/hub} hot loop, \texttt{shouldThrottle} is at
100.0 percent and \texttt{stampAndFanOut} is at 77.8 percent. The
lowest-coverage routine flagged is \texttt{ws/hub.route} at
50.0 percent (pending a refactor to expose its private branches
to unit tests).

\subsection{Fuzz Testing}
Go-native fuzz targets are committed in
\texttt{internal/frame/fuzz\_test.go} and
\texttt{internal/dispatcher/fuzz\_test.go}. They are run
opportunistically by developers on branch builds and do not gate CI.

%======================================================================
% 5. ROS2 Tests
%======================================================================
\section{ROS2 Tests}
\label{sec:ros2}

\subsection{Scope}
Fourteen test files split between C++ (ament\_cmake\_gtest) and
Python (ament\_python with \texttt{launch\_test}). Tests are
invoked by \texttt{colcon test --return-code-on-test-failure} in
\texttt{.github/workflows/ros2.yml}.

\begin{table}[H]
\centering
\caption{ROS2 test inventory.}
\label{tab:ros2-tests}
\small
\begin{tabularx}{\linewidth}{l X}
\toprule
\textbf{Package / Test} & \textbf{Verifies} \\
\midrule
star\_spi\_bridge / test\_spi\_driver.cpp          & SPI driver unit behavior \\
star\_spi\_bridge / test\_spi\_message\_converter  & Protobuf round-trip \\
star\_gateway\_bridge / test\_message\_converter   & Gateway bridge message codec \\
star\_safety\_monitor / test\_safety\_monitor      & Heartbeat, stall, obstacle logic \\
star\_perception / test\_hailo\_runner.py          & Hailo inference runner shim \\
star\_perception / test\_depth\_fusion.py          & Stereo + Hailo fusion \\
star\_bringup / launch\_test\_stereo\_topics.py    & Expected topics publish after launch \\
star\_simulation / test\_teleop\_timeout.py        & Teleop command timeout safety \\
star\_simulation / test\_ekf\_drift.py             & EKF drift bounds under simulated noise \\
star\_simulation / test\_stall\_detection.py       & Motor-stall triggers E-stop \\
star\_simulation / test\_slam\_stability.py        & SLAM convergence on synthetic scan \\
star\_simulation / test\_heartbeat\_loss.py        & Heartbeat-miss -> E-stop \\
star\_simulation / test\_nav2\_recovery.py         & Nav2 recovery behavior tree \\
star\_simulation / test\_obstacle\_estop\_latency.py & Obstacle -> E-stop latency budget \\
\bottomrule
\end{tabularx}
\end{table}

\subsection{Results}
All ROS2 tests pass on CI. The CI workflow uploads full test-result
artifacts from \texttt{star-ros2/build/**/test\_results/} on failure,
and applies \texttt{uncrustify} C++ formatting as a separate gate.

\subsection{Safety-Critical Tests}
The \texttt{star\_simulation} package contains four safety-critical
behavior tests. Each is a \texttt{launch\_test} that boots the
required ROS2 nodes in-process, injects the stimulus, and asserts the
expected response within a latency budget.

%======================================================================
% 6. Compliance Engine Tests (Python)
%======================================================================
\section{Compliance Engine Tests}
\label{sec:comp}

\subsection{Scope}
Twelve pytest modules in
\texttt{final\_capstone/compliance-engine/tests/} verify the
ADA-check algorithms and ROS2 node wrappers.

\begin{table}[H]
\centering
\caption{Compliance engine test modules.}
\label{tab:comp-tests}
\small
\begin{tabularx}{\linewidth}{l X}
\toprule
\textbf{Test module} & \textbf{Focus} \\
\midrule
test\_plane\_segmentation.py    & RANSAC plane extractor (ramp slope core) \\
test\_yolo\_fusion.py           & YOLO + LiDAR fusion glue \\
test\_doorway\_lidar\_detector.py & Door detection from scan \\
test\_obstacle\_clusterer.py    & Clustering of obstacle points \\
test\_cane\_zone\_filter.py     & Cane-sweep zone extraction \\
test\_corridor\_medial\_axis.py & Corridor medial-axis path width \\
test\_imu\_jolt\_detector.py    & IMU-based trip-hazard detection \\
test\_door\_offset\_calibration.py & Door offset calibration \\
test\_floor\_frame.py           & Ground-plane frame estimation \\
test\_jamb\_plane\_fitter.py    & Door jamb plane fit \\
test\_nodes\_smoke.py           & ROS2 node smoke tests \\
test\_nodes\_behavior.py        & Behavior tests on canned scans \\
\bottomrule
\end{tabularx}
\end{table}

\subsection{Results and Coverage}
Pytest is invoked with
\texttt{pytest --maxfail=1 --cov=star\_compliance
--cov-report=xml --cov-report=term --junitxml=pytest-report.xml}.
All tests pass at HEAD; coverage XML is uploaded as a workflow
artifact. The approval stage of the four-phase
\texttt{compliance-engine-ci.yml} workflow blocks merge if any unit
test fails or if the integration-phase smoke test does not complete.

%======================================================================
% 7. Protocol Buffers Tests
%======================================================================
\section{Protocol Buffers Tests}
\label{sec:proto}

\subsection{Cross-Language Serialization Tests}
Three parallel test harnesses validate that Go, TypeScript, and
nanopb encode identical bytes for the same canonical fixtures.

\begin{table}[H]
\centering
\caption{Protobuf serialization test suites.}
\label{tab:proto-tests}
\small
\begin{tabularx}{\linewidth}{l l X}
\toprule
\textbf{Language} & \textbf{Location} & \textbf{Coverage} \\
\midrule
Go        & \texttt{star-proto/tests/go/}         & 5 tests (telemetry, config,
                                                    firmware update, motor control,
                                                    serialization round-trip) \\
TypeScript & \texttt{star-proto/tests/typescript/} & 3 tests (telemetry,
                                                     serialization, motor control) \\
nanopb (C) & \texttt{star-proto/tests/nanopb/}    & Unity-based C serialization \\
\bottomrule
\end{tabularx}
\end{table}

\subsection{Lint and Breaking-Change Gates}
The \texttt{.github/workflows/proto.yml} workflow runs, in order:
\texttt{buf lint} (style), \texttt{buf format} (formatting),
\texttt{buf build} (schema integrity), and \texttt{buf breaking}
against main on pull requests. A breaking change blocks merge unless
the PR updates firmware, gateway, and ROS2 in the same change.

\subsection{Results}
All three serialization harnesses pass at HEAD. Buf lint, format,
and build are green; no open breaking-change violations.

%======================================================================
% 8. UI Tests
%======================================================================
\section{User Interface Tests}
\label{sec:ui}

\subsection{Scope}
The legacy React operator UI at \texttt{star-ui/} is the subject of
a single Vitest suite at
\texttt{star-ui/src/hooks/useGamepad.test.ts}. It verifies the
gamepad hook's axis mapping, dead-zone handling, and \qty{50}{\hertz}
command cadence. Coverage configuration is present but not gated in
CI, because the primary operator dashboard (Grafana and Lichtblick)
replaced the React UI as the main entry point late in development.

\subsection{Status}
\begin{table}[H]
\centering
\caption{UI test status.}
\small
\begin{tabular}{l l}
\toprule
\textbf{Suite} & \textbf{Status} \\
\midrule
\texttt{useGamepad.test.ts} & \pass{} \\
\bottomrule
\end{tabular}
\end{table}

The gap between this single suite and the thirty-panel inventory
of the original React UI is acknowledged explicitly: the UI is
\statuslabel{STRETCH} as delivered, since the Grafana and Lichtblick dashboards
became the primary surfaces the grader will exercise. Tests for
those dashboards exist at the protocol layer (gateway and
compliance-engine tests cover every endpoint they invoke).

%======================================================================
% 9. Integration and Simulation
%======================================================================
\section{Integration and Simulation Tests}
\label{sec:int}

\subsection{Virtual RX72N}
\texttt{star-gateway/cmd/virtual\_rx72n/} is an in-process Go program
that reimplements the RX72N SPI peripheral surface against the
gateway's transport layer. It runs inside the gateway test binary,
so the full protocol stack (framing, FEC, HARQ, gRPC) can be
exercised end-to-end without any wire. The associated test files
\texttt{main\_test.go} and \texttt{control\_frames\_test.go} ensure
the simulator itself behaves identically to the real firmware at
the observable layer.

\subsection{End-to-End HIL Harness}
\texttt{star-gateway/test/e2e/hil\_test.go} drives the real RX72N
over SPI from a test harness. It is skipped in standard \texttt{go
test} invocations and triggered only by
\texttt{.github/workflows/firmware-hil-nightly.yml} on a dedicated
HIL fixture. Results are published as GitHub Actions artifacts.

\subsection{ROS2 Launch Tests}
Four behaviorally critical flows are covered by
\texttt{launch\_test} harnesses:
\begin{itemize}[leftmargin=1.4em,itemsep=2pt]
\item Teleop timeout on missing heartbeat
      (\texttt{test\_teleop\_timeout.py})
\item EKF divergence bounds
      (\texttt{test\_ekf\_drift.py})
\item Motor stall -> E-stop sequence
      (\texttt{test\_stall\_detection.py})
\item Obstacle -> E-stop end-to-end latency
      (\texttt{test\_obstacle\_estop\_latency.py})
\end{itemize}

\subsection{Protocol Bridge Layer Sanity}
The gateway test suite exercises the complete protocol stack in
simulation: a gRPC call enters at \texttt{internal/service/},
flows through \texttt{internal/harq/}, \texttt{internal/fec/}, and
\texttt{internal/frame/}, exits via \texttt{internal/transport/}
into the virtual RX72N, is decoded back up the stack, and the
response is validated against the original request. This is the
closest thing STAR has to a bit-exact protocol regression harness.

%======================================================================
% 10. CI/CD Workflows
%======================================================================
\section{CI/CD Workflows}
\label{sec:ci}

\subsection{Workflow Inventory}
Fifteen GitHub Actions workflows live under
\texttt{.github/workflows/}. Table~\ref{tab:workflows} lists the
enforcement-bearing workflows.

\begin{table}[H]
\centering
\caption{Enforcement-bearing CI workflows.}
\label{tab:workflows}
\small
\begin{tabularx}{\linewidth}{l X}
\toprule
\textbf{Workflow} & \textbf{What it enforces} \\
\midrule
firmware-unit-tests.yml     & Unity suite, lcov HTML, 100/100/100 coverage
                              on libs/ \\
gateway.yml                 & Seven jobs: proto gen, build+test (race),
                              lint, security, integration, benchmark,
                              summary; \texttt{MIN\_COVERAGE=65} \\
proto.yml                   & \texttt{buf lint/format/build},
                              breaking-change detection against main;
                              Go, TS, and nanopb test runs \\
ros2.yml                    & \texttt{colcon test}, uncrustify format,
                              artifact upload on failure \\
compliance-engine-ci.yml    & Four phases: lint, unit, integration,
                              approval; pytest + coverage + junit \\
firmware-hil-nightly.yml    & Hardware-in-the-loop, nightly \\
ascii-check.yml             & 7-bit ASCII on every source file \\
copyright-check.yml         & Header presence / consistency \\
firmware-build-verify.yml   & RX72N release build succeeds \\
stm32-firmware.yml          & STM32 firmware build (out of scope here) \\
firmware-toolchain-image.yml & GNURX container image cache \\
stm32-toolchain-image.yml   & STM32 container image cache \\
proto-gen.yml               & Regenerated artifacts match committed files \\
since-version-check.yml     & \texttt{@since} tags present on public APIs \\
firmware.yml                & RX72N firmware top-level CI orchestration \\
\bottomrule
\end{tabularx}
\end{table}

\subsection{Pre-Commit Hooks}
\texttt{.pre-commit-config.yaml} installs buf v1.50.1 hooks for
\texttt{buf-lint}, \texttt{buf-format}, and \texttt{buf-generate}
against \texttt{star-proto/proto}. The repository additionally uses a
local hook at \texttt{scripts/git/pre-commit} to enforce the 7-bit
ASCII policy; the same check is duplicated in
\texttt{ascii-check.yml} so local hook bypass does not defeat CI.

\subsection{Commit-Gate Summary}
A pull request cannot merge to main unless: firmware coverage holds
at 100 percent on \texttt{libs/}; gateway coverage is at least 65
percent; all ROS2 colcon tests pass; all compliance-engine pytest
modules pass; buf lint and breaking-change checks are clean; 7-bit
ASCII is satisfied; copyright headers are present; and generated
protobuf artifacts match the committed files.

%======================================================================
% 11. Test Results Consolidation
%======================================================================
\section{Consolidated Test Results}
\label{sec:results}

Table~\ref{tab:consolidated} is the top-level scorecard of every
test category delivered on main at the date of this report.

\begin{table}[H]
\centering
\caption{Consolidated test scorecard.}
\label{tab:consolidated}
\small
\begin{tabularx}{\linewidth}{l l l X}
\toprule
\textbf{Category} & \textbf{Count} & \textbf{Status} & \textbf{Notes} \\
\midrule
Firmware unit (Unity)            & 60 ctest entries & \pass{} & 100\% coverage on libs/ \\
Firmware bring-up (bench)        & 8 dirs       & \pass{} & All on-bench \\
Manual PCB validation            & 9 checks     & \pass{} & Continuity, motor load, probing (Table~\ref{tab:pcbvalid}) \\
Gateway unit (Go)                & 46 files / 14 pkgs & \pass{} & 74\% coverage, 65\% gate \\
Gateway fuzz (Go)                & 2 targets    & advisory & Not gated in CI \\
ROS2 colcon                      & 14 tests     & \pass{} & All packages \\
ROS2 safety behavior             & 4 launch tests & \pass{} & E-stop, stall, EKF, teleop \\
Compliance pytest                & 12 modules   & \pass{} & Coverage XML uploaded \\
Protobuf Go                      & 5 tests      & \pass{} & Serialization round-trip \\
Protobuf TS                      & 3 tests      & \pass{} & Serialization round-trip \\
Protobuf nanopb (C)              & 1 harness    & \pass{} & Unity-based \\
UI Vitest                        & 1 suite      & \pass{} & Gamepad hook only \\
Integration (virtual\_rx72n)     & 2 test files & \pass{} & In-process simulator \\
HIL nightly                      & 1 harness    & \pass{} & GitHub Actions artifact \\
Pre-commit ASCII                 & enforced     & \pass{} & Local + CI \\
CI gates                         & 15 workflows & \pass{} & On HEAD \\
\bottomrule
\end{tabularx}
\end{table}

The single visible shortfall is the sparseness of the React UI test
suite. Because the operator surface migrated to Grafana and
Lichtblick mid-semester, the React UI's test coverage was not
expanded to match its original thirty-panel scope; instead, the
new dashboards rely on the gateway and compliance-engine test
suites to validate every endpoint they invoke. This is recorded in
Section~\ref{sec:ui}.

%======================================================================
% 12. Defects, Gaps, and Mitigations
%======================================================================
\section{Known Defects, Gaps, and Mitigations}
\label{sec:defects}

\subsection{Open Gaps}
\begin{itemize}[leftmargin=1.4em,itemsep=2pt]
\item \textbf{UI coverage} -- single Vitest suite against a
      thirty-panel inventory. Mitigation: the Grafana/Lichtblick
      dashboards are exercised at the API layer by gateway tests;
      the React UI is \statuslabel{STRETCH} rather than required.
\item \textbf{\texttt{ws/hub.route} at 50 percent} -- a private-branch
      coverage hole. Mitigation: a refactor is scheduled for
      post-capstone work; gateway total coverage still exceeds the
      65 percent gate
      with headroom.
\item \textbf{Fuzz targets advisory only} -- Go fuzzing is not
      wired into CI time budgets. Mitigation: developers run fuzz
      targets opportunistically on frame and dispatcher changes.
\item \textbf{Online PID autotuning} -- not tested because not
      implemented (see SDD, Section on Known Limitations). PID
      gains are produced offline in MATLAB and verified via the
      \texttt{rx\_pid} unit test against a reference trajectory.
\end{itemize}

\subsection{Closed-Loop HIL Validation}
The obstacle-to-E-stop latency test
(\texttt{test\_obstacle\_estop\_latency.py}) is the most
safety-critical test in the suite. It verifies that a synthetic
obstacle appearing inside the \qty{10}{\centi\meter}
\texttt{obstacle\_estop\_distance} threshold triggers a
\texttt{/star/estop = true} publication inside one diagnostic
period (\qty{100}{\milli\second}), and that the gated velocity
output returns to zero immediately thereafter. The test is green
at HEAD.

\subsection{Regression Policy}
Breaking a CI gate requires the responsible PR to be reverted or
the gate itself to be explicitly relaxed in the same PR. The
project has no "skip CI" escape hatch; both branch protection and
required status checks are enabled on main.

%======================================================================
% 13. Conclusions
%======================================================================
\section{Conclusions}
\label{sec:conclusions}

The STAR software delivers a broad, well-instrumented test suite:
60 firmware ctest entries at 100 percent library coverage (the
e2-studio JUnit snapshot from 2026-03-01 captured 46 of these
suites), 46 Go test files across 14 packages at 74 percent gateway
coverage, fourteen ROS2 tests including four safety-critical
launch tests, twelve compliance-engine pytest modules, nine
Protobuf serialization tests spanning three languages, eight
on-bench hardware bring-up programs, nine manual PCB validation
checks (power-rail continuity, motor free-spin and loaded-chassis
load tests with current measurements, scope-probed PWM and
encoder waveforms, AD2-probed I2C/UART/SPI lines), and fifteen
CI workflows that enforce the gates described in
Section~\ref{sec:ci}. Every merge to main passes the full suite.

At the date of this report, the only visible gap is the React UI
suite, which is acknowledged as \statuslabel{STRETCH} because the dashboarding
surface migrated to Grafana and Lichtblick. Every other subsystem
is at or above its intended coverage bar.

\section*{Appendix A: Reproducing the Test Runs}
\addcontentsline{toc}{section}{Appendix A: Reproducing the Test Runs}

\begin{lstlisting}[style=shellcmd]
# Firmware unit tests
cd star-rx72n-firmware/tests
cmake -B build && cmake --build build
ctest --test-dir build --output-on-failure

# Gateway tests
cd star-gateway
go test -race -coverprofile=coverage.out -covermode=atomic ./...
go tool cover -func=coverage.out | tail -n 1   # total %

# ROS2 tests
cd star-ros2
colcon build --symlink-install
colcon test --return-code-on-test-failure --event-handlers console_direct+
colcon test-result --all

# Compliance engine
cd final_capstone/compliance-engine
pytest --cov=star_compliance --cov-report=term tests/

# Protobuf
cd star-proto
buf lint && buf format --diff && buf build
( cd tests/go        && go test -v ./... )
( cd tests/typescript && npm test )

# UI
cd star-ui
npm install && npm test
\end{lstlisting}

\section*{Appendix B: Key Test Artifacts}
\addcontentsline{toc}{section}{Appendix B: Key Test Artifacts}

\begin{itemize}[leftmargin=1.4em,itemsep=2pt]
\item \texttt{e2-studio-star-rx72n-firmware/tests/build/test-results.xml}
      -- firmware JUnit XML snapshot, 46 of 60 suites, 0 failures
      (CI runs the full 60-suite ctest tree).
\item \texttt{star-gateway/coverage.out},
      \texttt{star-gateway/coverage.html},
      \texttt{star-gateway/coverage\_transport.out}
      -- gateway coverage outputs.
\item \texttt{.github/workflows/*.yml}
      -- all CI workflow definitions.
\item \texttt{.pre-commit-config.yaml}
      -- committed pre-commit hooks.
\item \texttt{scripts/git/pre-commit}
      -- ASCII enforcement hook.
\end{itemize}

\end{document}
